\documentclass[11pt,a4paper]{article}
\usepackage{amsmath,amssymb,amsthm}
\usepackage[margin=2.5cm]{geometry}
\usepackage{booktabs}
\usepackage{enumitem}
\usepackage{xcolor}
\usepackage[
    colorlinks=true,
    linkcolor=blue,
    citecolor=darkgray,
    urlcolor=blue,
    bookmarks=true,
    bookmarksnumbered=true,
    unicode=true
]{hyperref}

\widowpenalty=10000
\clubpenalty=10000

\newcommand{\Z}{\mathbb{Z}}
\newcommand{\Q}{\mathbb{Q}}
\newcommand{\C}{\mathbb{C}}
\newcommand{\R}{\mathbb{R}}
\newcommand{\BK}{\mathrm{BK}}
\newcommand{\BU}{\mathrm{BU}}
\newcommand{\Fl}{\mathrm{Fl}}
\newcommand{\BKg}{B\kappa}
\newcommand{\tors}{\,\mathrm{tors}}

\begin{document}

\begin{center}
{\large\bfseries A Referee-Style Read of Proposition 6.16}\\[4pt]
{\itshape The cup square on the flag quotient: a hand derivation}\\[6pt]
{\small Instruments manuscript v10, source lines 1857--2084,}\\
{\small with the integration points: Lemma 6.17, Remarks 6.18--6.19, the open-problem list,}\\
{\small and the abstract/intro/conclusion provenance phrases}\\[6pt]
{\small Wave 30 \,\textbullet\, 2026-09-13 \,\textbullet\, channel-supp-augmented}
\end{center}
\vspace{4pt}

\noindent\textbf{Charge and method.} Every step of the seven-step proof was
re-derived by hand, independently of the manuscript and of the wave-13C/28
machine data, and every checkable arithmetic claim was re-computed by an
exact-integer script (\texttt{w30\_referee\_check.py}) written from scratch
for this read. The machine tuple was consulted only at the end, as a
comparison, never as an input --- the same discipline the wave itself claims,
and the claim survives this referee's audit of it.

\section*{1. Verdict, up front}

\textbf{The Proposition is correct in its principal assertions and the proof
route is sound.} The assertion $x^2\ne 0$ in $H^4(B;\Z)$ --- the decision bit
that discharges the machine premise of Lemma~6.17 --- is proved by a valid
chain (the homogeneous-space model, the pullback of the universal fibration,
Borel's transgression with the $\varepsilon=2$ pin, the collapsed
$BT^3$-page, the torsion-survival crux, and the edge identification); the
additive tuple
$(\Z,0,\Z/3,\Z/3,\Z/3,\Z/3,\Z)$ is correct; and the claims that $x$ generates
$H^2$ and $x^2$ generates $H^4$ are correct.

\textbf{One substantive correction is required (Finding~R1):} the final
clause of the statement --- ``and a discriminant class generating $H^6$'' ---
is false as stated. In $H^6(B;\Z)\cong\Z$ the discriminant class is
\emph{twice} a generator. The error originates in Step~2's integral module
statement, which is true rationally but false integrally (an index-2
sublattice, with the index growing in higher degrees). The correction is
local: the groups, the crux, the lemma, and the downstream theorem chain
(Theorem~6.13's plane-valued extension, \texttt{thm:state-d2}, and
$\delta_2(\mathcal D(\C^3))=4/3$) are all unaffected.

\textbf{Recommendation: accept with a minor revision} --- R1's local
correction plus the six one-line repairs R2--R7 below. In journal terms:
\emph{correct and resubmit; no new mathematics needed.}

\section*{2. What was independently re-verified (all confirmed)}

\begin{itemize}[leftmargin=1.4em,itemsep=2pt]

\item \textbf{Step 0 (the homogeneous-space model).}
$\Fl=U(3)/T^3$ with the flag lines the column spans; the deck transformation
$c$ is right multiplication by the permutation matrix $\sigma$ (column
shift), which normalises $T^3$; the $\langle c\rangle$-orbits are exactly the
right cosets of $K=T^3\rtimes\langle\sigma\rangle$, so $B=U(3)/K$. Verified,
including the harmless convention ambiguity ($\sigma$ versus $\sigma^{-1}$
generates the same $K$).

\item \textbf{Step 1 (the fibration and its transgressions).} The Borel
construction $EU(3)\times_K U(3)\to\BK$ is the associated bundle with fibre
the left $K$-space $U(3)$; the total space is weakly equivalent to
$K\backslash U(3)\cong B$; the map $[e,g]\mapsto eg$ exhibits it as the
pullback of the universal $U(3)$-bundle along $B\iota$, with the identity on
fibres --- so the Serre-SS differentials are the pulled-back universal ones by
naturality, and the local system is trivial because left translations by
elements of the connected $U(3)$ are homotopic to the identity. The
transgression $d_{2i}(z_{2i-1})=c_i$ is classical (Borel), as cited.

\item \textbf{Step 2 (the cohomology of $\BK$) --- the groups.} The collapse
argument of the $BT^3\to\BK\to BC_3$ page is airtight: with the
permutation-orbit decomposition of each $\operatorname{Sym}^d$, an even-$r$
differential targets an odd fibre degree (where $H^{\mathrm{odd}}(BT^3)=0$)
and an odd-$r$ differential targets an odd, positive group degree (permutation
modules are acyclic; the trivial summands have $H^{\mathrm{odd}}(C_3;\Z)=0$).
The free ranks equal the Burnside orbit counts (re-computed for $d\le8$:
$1,1,2,4,5,7,10,12,15$), and the torsion sits in bidegrees $(2k,6a)$, $k\ge1$,
one $\Z/3$ each. Confirmed.

\item \textbf{Step 3 (the torsion product rule) --- the arithmetic.} All the
$\tau$-products used in the proof were re-computed from the norm model
$H^2(C_3;M_d)=M_d^{C_3}/NM_d$: $\tau\sigma_1=\tau\sigma_2=\tau\Delta=
\tau\sigma_1^3=\tau\sigma_1\sigma_2=0$ (the $(\chi_1\chi_2\chi_3)$-coefficients
$0,0,0,6,3$ respectively) and $\tau\sigma_3\ne0$ (coefficient $1$),
generating the $(2,6)$-slot; $\tau^k\ne0$ for all $k$. The display
$H^4(\BK)\cong\Z\sigma_1^2\oplus\Z\sigma_2\oplus\Z/3\cdot\tau^2$ is correct
\emph{as a lattice}: the degree-2 orbit sums
$\chi_1^2+\chi_2^2+\chi_3^2=\sigma_1^2-2\sigma_2$ and
$\chi_1\chi_2+\chi_2\chi_3+\chi_3\chi_1=\sigma_2$ generate exactly
$\langle\sigma_1^2,\sigma_2\rangle$ (index~1, verified). This matters: the
degrees the crux needs are exactly the safe ones.

\item \textbf{Step 4 (the transgression values).} The Chern-root restriction
$c_i(\chi)=\sigma_i(\chi)$ on the fibre; the splitting $s\colon BC_3\to\BK$
through $\sigma\mapsto(1;\sigma)$; the composite classifies the regular
representation $1\oplus\omega\oplus\omega^2$, whose Chern classes in
$\Z[u]/(3u)$ are $c_1=3u=0$, $c_2=2u^2$, $c_3=0$ (re-computed:
$(1+0u)(1+1u)(1+2u)=1+3u+2u^2$). The section kills positive fibre degree, so
$2u^2=\varepsilon u^2$ pins \textbf{$\varepsilon=2$} and $0=\varepsilon'u^3$
pins \textbf{$\varepsilon'=0$}; $s^*(\tau^2)=u^2$ is the ring-map image. All
confirmed.

\item \textbf{Step 5 (the crux) --- the heart of the proof.} The incoming
images into the $(4,0)$-slot are $\langle\sigma_1^2\rangle$ (from
$d_2(\beta z_1)=\beta\sigma_1$; note $\tau\sigma_1=0$ kills the torsion part
of the image) and $\langle\sigma_2+2\tau^2\rangle$ (from $d_4(z_3)$); nothing
leaves $(4,0)$; the $(0,4)$-term dies under $d_2(z_1z_3)=\sigma_1z_3\ne0$
(injective, free target); the $(2,2)$- and odd slots vanish. Hence
\[
E_\infty^{4,0}=\bigl(\Z\sigma_2\oplus\Z/3\cdot\tau^2\bigr)
\big/\langle\sigma_2+2\tau^2\rangle\cong\Z/3,
\]
the Smith normal form of
$\bigl(\begin{smallmatrix}1&2\\0&3\end{smallmatrix}\bigr)$
giving $[\tau^2]$ of order $3$, nonzero, generating
(with $[\sigma_2]=-2[\tau^2]$). \textbf{The crux is correct.} (See~R5 on the
phrasing.)

\item \textbf{Step 6 (the edge).} The edge of the fibration's spectral
sequence is the projection pullback; the cover $\Fl\to B$ is the
$T^3$-quotient of the principal $K$-bundle $U(3)\to B$, hence classified by
$\pi\circ\gamma$ with $\gamma$ the classifying map --- functoriality of the
associated-bundle construction, rigorous. So $\kappa=\pi\circ\gamma$ and
$x=\kappa^*(u)=\gamma^*(\tau)$ (see~R6 for a convention remark), and
$x^2=\gamma^*(\tau^2)=[\tau^2]\ne0$. Confirmed --- \textbf{the main assertion
holds.}

\item \textbf{Step 7 (the remaining degrees).} Degree~2:
$H^2(B)=\Z/3$, generated by the edge image of $\tau$ (that is, $x$).
Degree~3: $d_2(\sigma_1z_1)=\sigma_1^2\ne0$ kills the free part of $(2,1)$
while $d_2(\tau z_1)=\tau\sigma_1=0$ leaves $\Z/3\cdot\tau z_1$; $z_3$ dies
under $d_4$ --- and the image $\sigma_2+2\tau^2$ has \emph{infinite order},
so $d_4$ is injective on $\Z z_3$; this is exactly the point where the
derivation correctly diverges from the lens model (there the transgression
image is torsion and $3z_3$ survives). Degree~5: the free $(4,1)$-classes die
under $d_2$ ($\tau^2\sigma_1=\tau(\tau\sigma_1)=0$ survives); $\tau z_3$ dies
under $d_4(\tau z_3)=\tau(\sigma_2+2\tau^2)=2\tau^3\ne0$; $z_5$ dies under
$d_6(z_5)=\sigma_3\ne0$; so $H^5=\Z/3\cdot\tau^2z_1$. Degree~6: the group is
$\Z$ (see R1 for the generator). The mixed terms die transversally as
claimed. The degree runs are all correct.

\item \textbf{The tuple.} $(\Z,0,\Z/3,\Z/3,\Z/3,\Z/3,\Z)$, independently,
and the UCT inversion matches the wave-13C/28 machine homology tuple
$(\Z,\Z/3,\Z/3,\Z/3,\Z/3,0,\Z)$ --- read only as a comparison, as the wave
states. Euler characteristic $2=6/3$, consistent.

\item \textbf{The Lemma 6.17 decision logic.} $x^2=\kappa^*(u^2)$ is the
edge image of $u^2$ through $E_\infty^{4,0}=\operatorname{coker}
d_3^{1,2}$; a nonzero image forces the cokernel nonzero; a
$\Z/3\to\Z/3$ map with nonzero cokernel is zero; so $d_3^{1,2}=0$ and
$H^3=H^4=\Z/3$ with $x^2$ generating. Sound; the page-level computations
($E_2^{1,2}=\Z/3$ via the norm/invariants on $H^2(\Fl)$, the vanishing
fixed-vector argument on $H^4(\Fl)$) re-verified.

\item \textbf{The lens-model validation (G7).} Re-derived by hand:
$SU(2)\to ESU(2)\times_{C_3}SU(2)\to BC_3$ with $d_4(z_3)=2u^2$ (the pin for
the two-dimensional representation $\omega\oplus\omega^2$) reproduces
$H^*(L(3;1))=(\Z,0,\Z/3,\Z)$ with the surviving $3z_3$ --- the mechanism,
and the contrast with the flag case, are exactly as the wave reports. The
$\R\mathrm{P}^2\times S^2$ companion case pins the other world ($x^2=0$,
$d_3$ an isomorphism), as recorded in Wave~24.
\end{itemize}

\section*{3. Finding R1 (substantive) --- the integral module statement and
the $H^6$-generator clause}

\textbf{What is claimed.} Step~2: ``the free part is
$\Z[\sigma_1,\sigma_2,\sigma_3]\oplus\Delta\cdot\Z[\sigma_1,\sigma_2,
\sigma_3]$'', justified by rank reproduction (the orbit count equals the
monomial count of the module). Step~3 displays
$H^6(\BK)\cong\Z\{\sigma_1^3,\sigma_1\sigma_2,\sigma_3,\Delta\}\oplus
\Z/3\cdot\tau^3$. Step~7 concludes $H^6(B)\cong\Z\cdot\langle\Delta\rangle$,
``the top class being the discriminant'', and the \emph{statement} asserts
``a discriminant class generating $H^6$''.

\textbf{What is wrong.} The rank-reproduction argument proves only that the
displayed module is a \emph{finite-index} sublattice of the invariant
(orbit-sum) lattice $M_d^{C_3}=H^{2d}(\BK)/\tors$ (the edge
$H^{2d}(\BK)\to E_\infty^{0,2d}=H^0(C_3;\operatorname{Sym}^d)$ is surjective
with torsion kernel). Integral equality fails from degree~3 on. The index,
computed exactly:
\[
\bigl[M_d^{C_3}:\text{module}\bigr]_{d=0,\dots,8}
=\bigl(1,\;1,\;1,\;2,\;2,\;4,\;8,\;16,\;32\bigr).
\]
The degree-3 witness: the orbit sum
$S_1=\chi_1^2\chi_2+\chi_2^2\chi_3+\chi_3^2\chi_1$ is invariant, and
\[
S_1+S_2=\sigma_1\sigma_2-3\sigma_3,\qquad
S_1-S_2=\Delta\qquad\Longrightarrow\qquad
2S_1=\sigma_1\sigma_2-3\sigma_3+\Delta,
\]
so the linear system $S_1=a\sigma_1^3+b\sigma_1\sigma_2+c\sigma_3+e\Delta$
forces $b=e=\tfrac12$ --- no integer solution. Over $\Q$ the module statement
is of course true (2 invertible), which is why the rational cohomology
remark in Step~7 survives.

\textbf{How it propagates.} Two paths, both confined to degree~6. First, the
Step~3 display is correct as an abstract group ($\Z^4\oplus\Z/3$) but false
as an identification via the named classes: they span an index-2 sublattice
of the free part; the true free part is the orbit-sum lattice
$\langle A_1,S_1,S_2,\sigma_3\rangle$ with
$A_1=\chi_1^3+\chi_2^3+\chi_3^3$. Second, Step~7's degree-6 quotient, run on
the true lattice $M_3^{C_3}/\langle\sigma_1^3,\sigma_1\sigma_2,\sigma_3
\rangle$, is $\cong\Z$ (free rank~1, no torsion --- verified), but the
generator is the class of the \textbf{orbit sum} $S_1$: the relations give
$[S_2]=-[S_1]$, $[A_1]=0$, $[\sigma_3]=0$, hence
\[
[\Delta]=[S_1-S_2]=2\,[S_1],
\]
\textbf{the discriminant class is twice a generator}. Verified exactly:
$\Delta-2S_1$ lies in the image lattice, while $\Delta-S_1$ and $S_1$ do not.
So both ``$\cong\Z\cdot\langle\Delta\rangle$'' (Step~7) and ``a discriminant
class generating $H^6$'' (the statement) are incorrect as stated; the group
$\Z$ is correct.

\textbf{What is not affected.} The crux and everything downstream of it: the
crux lives in fibre-degrees $\le2$, where the module statement is exact
(index~1, verified); the $\tau$-rule is a monomial computation, independent
of the module presentation; the classes $\sigma_1,\sigma_2,\sigma_3$ are
pinned as canonical classes (Chern-class pullbacks) in Step~4, not via the
module statement; and Lemma~6.17, the plane-valued theorem,
\texttt{thm:state-d2}, and $\delta_2(\mathcal D(\C^3))=4/3$ use only
$H^2/H^3/H^4$ and the decision bit. The additive tuple --- the object the
three independent routes agree on --- is untouched. The rational statement
$H^*(B;\Q)=\Q\oplus\Q\cdot\bar\Delta$ remains correct (and is the right
framing for the discriminant).

\textbf{Why the wave's gates did not catch it.} Gate~G3 of
\texttt{w29\_derivation\_check.py} compares the \emph{ranks} (orbit count
versus $\mathrm{nmons}(d)+\mathrm{nmons}(d-3)$) --- an index-2 sublattice is
invisible to a rank check. (G10 compares groups; $\Z$ matches $\Z$.) This is
a gate-design gap, not an execution failure: the gates certified exactly what
they tested, and the manuscript then asserted slightly more than the gates
certified.

\textbf{The repair (mechanical, about six lines of source).}
(i)~Step~2: replace the module equality by the orbit-sum lattice statement
(state the module over $\Q$ with the integral free ranks given separately;
optionally note the index). (ii)~Step~3: replace the $H^6(\BK)$ generating
set by the orbit sums $\{A_1,S_1,S_2,\sigma_3\}$, or display the group
abstractly. (iii)~Step~7: conclude $H^6(B;\Z)\cong\Z$ generated by the
class of the orbit sum $S_1$, with the discriminant twice a generator.
(iv)~The statement: replace ``a discriminant class generating $H^6$'' by
``$H^6(B;\Z)\cong\Z$, in which the discriminant class is twice a generator''
(or drop the clause --- the tuple is the operative content).
(v)~Add a lattice-index gate to the verifier (membership of $S_1$ in the
module lattice, or the index table above) so the correction is certified.
(vi)~The wave-29 note and the commit message record ``$H^6=\Z\cdot\Delta$'';
the correction wave should note the supersession.

\section*{4. Findings R2--R7 (minor; one-line repairs each)}

\begin{itemize}[leftmargin=1.4em,itemsep=3pt]
\item \textbf{R2 (Step 1).} The parenthetical justifying $d_r=0$ for
$r\notin\{2,4,6\}$ reads ``$r$ even forces an odd $q$-degree in the target,
where $H^{\mathrm{odd}}(U(3))=0$'' --- but $H^{\mathrm{odd}}(U(3))\ne0$
($z_1,z_3,z_5$ are odd-degree classes). The claim is true; the correct
justification is two-step: (i)~$r$ odd forces an odd, positive $p$-degree in
the target, where $H^{\mathrm{odd}}(\BK)=0$; (ii)~$r$ even, $r\ge8$ forces
the only possible source fibre-degrees to be $8$ and $9$ ($z_3z_5$,
$z_1z_3z_5$), and those classes are already killed by $d_2/d_4$ (e.g.\
$d_4(\beta z_3z_5)=\beta(\sigma_2+2\tau^2)z_5\ne0$ for every $\beta$, free or
torsion). Reword the parenthetical.

\item \textbf{R3 (Step 5).} The enumeration of the total-degree-4 terms
silently omits $(2,2)$. It vanishes because $H^2(U(3);\Z)=0$ --- one clause.

\item \textbf{R4 (Step 4).} The parenthetical ``the free generator is the
unique class restricting to $\sigma_1(\chi)$, since the torsion restricts to
zero'' does not establish uniqueness: the classes restricting to
$\sigma_1(\chi)$ are the coset $\sigma_1+b\tau$, $b\in\{0,1,2\}$. The
constant $b$ \emph{is} pinned --- by the same splitting the paragraph then
uses for $\varepsilon,\varepsilon'$: $s^*\iota^*c_1=c_1(\mathrm{reg})=3u=0$
and $s^*(\sigma_1+b\tau)=bu$ force $b=0$. Add the $c_1$ case to the pinning
sentence.

\item \textbf{R5 (Step 5).} ``A finite-order element never lies in a subgroup
generated by an infinite-order element'' rules out membership in each image
subgroup \emph{separately}, but the quotient is by their \emph{sum}. The gap
is closed by the displayed two-step quotient (first split off the
$\Z\sigma_1^2$ direct summand --- the image of $\tau^2$ is unchanged by a
direct-summand quotient --- and only then mod out by the single remaining
relation); the conclusion is correct (SNF-verified). One sentence would make
it airtight.

\item \textbf{R6 (Step 6).} The identification of the transported projection
with the classifying map $\gamma$ --- hence the exact equality
$x=\gamma^*(\tau)$ --- is correct in substance but depends on the alignment
of the deck-action convention with the $C_3$-structure on the associated
bundle. The robustness note: the only possible discrepancy is a sign in
$\Z/3$ ($x=\pm\gamma^*(\tau)$), and the conclusion $x^2=\gamma^*(\tau^2)$ is
invariant under it ($2^2=4\equiv1\bmod3$). A one-line remark removes the
dependence on the reader's convention bookkeeping.

\item \textbf{R7 (Step 3).} The claim $\tau\Delta=0$ is sensitive to the
choice of the $\Delta$-lift's torsion component (the $\tau$-rule controls
the leading graded component; a torsion tail of the lift could contribute a
deeper $\tau\cdot(\text{tail})$ term), and no pin is given for that tail. The
claim is never used downstream (the crux and the degree runs use only
$\tau\sigma_1=\tau\sigma_2=\tau^2\sigma_1=0$ and $\tau\sigma_3,\tau^3\ne0$,
all sound because the $\sigma_i$ are pinned pure by Step~4's splitting), so
nothing fails --- but either pin the tail, or mark the claim as
lift-dependent. (Also $\tau\sigma_1\sigma_2=0$ holds --- coefficient $3$ ---
if wanted for symmetry of the display.)
\end{itemize}

\section*{5. The integration points (assessed)}

\begin{itemize}[leftmargin=1.4em,itemsep=3pt]
\item \textbf{Lemma 6.17.} The rewritten proof --- the page-level computation,
the two-worlds statement, and the hand decision via Proposition~6.16 with the
machine retained as the cross-check --- is sound and correctly scoped. No
changes needed beyond R1's knock-on (none reaches the lemma: it uses degrees
$\le4$ only).
\item \textbf{Remark 6.19 (the machine certificate).} The
three-independent-routes framing (hand derivation, cellulation,
re-implementation) is fair and survives this read. Two phrases deserve
adjustment after R1: ``a self-contained invariant-theory step'' is where the
false module statement lives, and ``the invariant lattice \dots\ verified by
an exact-integer script'' overstates what G3 tested (ranks). Suggested
rewording: the invariant-theory input is the orbit decomposition and the
$\tau$-rule (monomial arithmetic, script-verified); the module presentation
is corrected per R1.
\item \textbf{Remark 6.18 (the flag-quotient literature).} The
differentiation from the Guerra--Jana programme (symmetric-group quotients
with field coefficients, in the service of Auerbach counting, versus the
cyclic intermediate quotient with integral coefficients and torsion
exponents pinned, in the service of the width chain) is accurate, and the
characterisation of the hand derivation as ``the integral lift of the
transgression route'' is now referee-confirmed in substance: the fibration
comparison, the Chern-class transgressions, and the invariant-theoretic
replacement of the field-coefficient module algebra are all present and
correct; the torsion-survival argument (Step~5) is exactly the layer the
field algorithms do not address. No changes needed.
\item \textbf{Open problem (vi).} The discharge is substantively valid: what
the open problem asked for (the hand derivation of the cohomology deciding
the lemma) is delivered. The R1 correction does not reopen it; the remaining
neighbourhood direction stated there (the unstable integral cohomology of
the symmetric-group quotients) is correctly identified.
\item \textbf{Abstract/intro/conclusion provenance phrases}
(``hand-derived, machine-certified, and independently re-implemented''):
accurate for what they assert (the additive cohomology of $B$ and the
cup-square decision); the main article's sync clauses are likewise fine. No
changes required by this read beyond R1's local fixes in the instruments
paper.
\end{itemize}

\section*{6. The referee's summary of the mathematical content}

The route is a genuine integral lift of the Guerra--Jana transgression
mechanism, not a notational variant of it: the pullback identification of the
Borel fibration (Step~1) is what makes the transgressions \emph{known}
rather than re-computed; the collapse of the $BT^3$-page (Step~2) is the
replacement for the field-coefficient module algebra; the $\varepsilon=2$ pin
(Step~4) is a clean use of the splitting section and the regular
representation; and the torsion-survival crux (Step~5) --- three lines of
order arithmetic that decide an element's survival against infinite-order
relations --- is the integral layer the cited programme does not touch. The
proof is checkable line by line, which was the point of the wave, and this
referee's independent re-computation confirms it in every load-bearing
place. The single defect found (R1) sits in the decorative layer --- the
ring/module presentation and the top-degree generator --- which is also
where a machine cross-check is least able to help: torsion-free rank-one
groups do not expose generator ambiguities. That is a fitting place for the
residual human error, and a fitting note on which to recommend revision
rather than rejection.

\section*{7. Artifacts of this wave}

\begin{itemize}[leftmargin=1.4em,itemsep=2pt]
\item \texttt{w30\_referee\_check.py} --- the independent exact-integer
verifier (seven check families: the module indices per degree; the degree-3
witness; the $\tau$-rule versus the norm computations; the
regular-representation Chern classes; the crux SNF and the order of
$[\tau^2]$; the degree-6 quotient with the $[\Delta]=2[S_1]$ membership
tests; the free-rank inventory).
\item \texttt{w30\_referee\_check\_output.txt} --- the transcript (all checks
pass in the directions asserted; the two \emph{disconfirming} results are the
module index table and the $[\Delta]=2[S_1]$ triple membership test, i.e.\
exactly the R1 evidence).
\item The full note \texttt{WAVE30\_REFEREE\_READ.md} in the \texttt{glm/}
folder of the companion repository, mirrored with the script and transcript;
committed and pushed with the session token (handled via a transient
\texttt{GIT\_ASKPASS} helper, never echoed, never committed).
\end{itemize}

\end{document}
